\documentclass[11pt]{article}
\usepackage[utf8]{inputenc}
\usepackage{geometry}
\geometry{a4paper, margin=1in}
\usepackage{titlesec}
\usepackage{enumitem}
\usepackage{xcolor}

\title{\textbf{Strategic Diagnostic Snapshot™}}
\author{RankPilot Intelligence}
\date{\today}

\begin{document}

\maketitle

\section*{I. Executive Summary}
\textbf{Firm:} González de Araujo Consultores \\
\textbf{Practice Area:} Data Protection \\
\textbf{Detected Practice Model:} Standard

\section*{II. Structural Positioning}
\textbf{Current Tier Assessment:} Unranked / New Entry \\
\textbf{Strategic Confidence:} 100.0\% \\
\textbf{Advantage:} Standard Market Offerings

\section*{III. Portfolio Analysis (Key Matters)}
\begin{itemize}
\item \textbf{Grupo Hermes Data Protection Framework} (Client: Grupo Hermes | Value: N/A)\\ The significance of the mandate for Grupo Hermes, a Mexican diversified infrastructure and services conglomerate, lies in its translation of data-protection requirements into defined governance processes with a measurable operational result. Rather than treating privacy as a purely documentary exercise, the firm was engaged to design and implement a thorough data protection and governance framework, placing the management of personal data flows and rights requests at the centre of the client’s internal procedures.

For Grupo Hermes, developing privacy notices, implementing policies, and establishing ARCO rights management procedures formed core and distinct components of the firm’s central role in the engagement. The team mapped personal data flows to identify the processing activity that the framework needed to address. On that foundation, it developed privacy notices, implemented policies, and established ARCO rights management procedures, linking the design of the governance framework to practical mechanisms for the handling of ARCO requests.

Arturo González de Araujo and Ernesto Zavala led the matter. Their work covered both the design and implementation of the thorough data protection and governance framework and the specific procedures required for ARCO rights management. The outcome was concrete: Grupo Hermes achieved 100% compliance in handling ARCO requests and avoided regulatory sanctions. This result gives the matter particular evidential weight as an example of a framework that was not only designed and implemented, but also applied through defined procedures capable of delivering full compliance in the handling of ARCO requests.
\item \textbf{Mega Direct Data Protection Advisory} (Client: MEGA DIRECT, S.A. de C.V. | Value: N/A)\\ Protection of MEGA DIRECT, S.A. de C.V., Customer experience, call center, and marketing services provider's databases and the continuity of its business operations sit at the centre of this significant data-protection mandate. The engagement demonstrates the team’s ability to represent a corporate client where the legal treatment of information assets is directly connected to ongoing operations in a highly regulated environment.

MEGA DIRECT, S.A. de C.V. instructed the firm on complex data protection and database-related matters. The work has included representing the client in legal proceedings concerning the acquisition, processing, storage and protection of data. Each of these stages is expressly material to the mandate: data acquisition, processing, storage and protection form the factual and legal focus of the proceedings in which the team has acted for MEGA DIRECT, S.A. de C.V.

Arturo González de Araujo, Ernesto Zavala and Damián Ortega lead the representation. Their involvement places the matter in a contentious setting, rather than limiting the engagement to documentation or internal policy work, and requires the team to address the legal issues arising from the client’s databases while preserving the continuity of its business activities. The concrete outcome identified by the submission is the maintenance of business continuity for MEGA DIRECT, S.A. de C.V. while the firm represents it in proceedings involving regulated data acquisition, processing, storage and protection.
\item \textbf{Biocodex Data Protection Framework} (Client: Biocodex | Value: N/A)\\ Biocodex, a global pharmaceutical company focused on probiotics and specialty healthcare, has a Data Protection Framework mandate that is significant because it placed responsibility for personal data protection within a dedicated institutional function, rather than treating the exercise solely as a set of standalone documents. Biocodex instructed the team in connection with the design and implementation of a structured personal data protection framework, an engagement directed at organising how personal-data obligations would be addressed within the client’s operations.

Arturo González de Araujo, Ernesto Zavala and Damián Ortega led the advice on designing and implementing that structured framework. Their work supported the establishment of a Personal Data Protection Department, creating a defined internal department through which the client could address personal-data protection requirements. The mandate therefore combined framework design with the implementation of an institutional structure responsible for the subject matter.

The work was expressly directed toward ensuring compliance with Mexican regulations while protecting sensitive data and trade secrets. The establishment of the Personal Data Protection Department is the concrete institutional outcome recorded for Biocodex: it provided a dedicated internal function connected to the structured personal data protection framework and to the handling of sensitive data and trade secrets under Mexican regulatory requirements.
\item \textbf{Hotel Riazor Data Protection Enhancement} (Client: Hotel Riazor | Value: N/A)\\ The significance of Hotel Riazor, Business hotel offering lodging and event services in Mexico City, Data Protection Enhancement lies in its focus on embedding personal-data governance within the client’s day-to-day operational and compliance arrangements, rather than addressing personal data protection as a standalone policy exercise. The mandate provided a concrete basis for reinforcing the internal structures responsible for lawful data handling.

Hotel Riazor, integrated into broader operational efficiency and compliance programmes, was supported by the firm in strengthening its personal data protection and governance framework. The engagement centred on reinforcing Hotel Riazor’s Personal Data Department and implementing policies designed to ensure lawful data handling. By linking the Personal Data Department’s reinforcement with the implementation of those policies, the work addressed both the client’s internal governance function and the operational application of lawful data-handling requirements.

Arturo González de Araujo, Ernesto Zavala and Damián Ortega led the mandate. Their role encompassed supporting the reinforcement of Hotel Riazor’s Personal Data Department, strengthening the client’s personal data protection and governance framework, and implementing policies to ensure lawful data handling. The resulting institutional measures comprised a reinforced Personal Data Department and implemented data-handling policies, integrated into broader operational efficiency and compliance programmes.
\item \textbf{Grupo Excelsior Data Protection and Compliance} (Client: Grupo Excelsior | Value: N/A)\\ Regularisation and restructuring of operations formed the central significance of the Grupo Excelsior, one of Mexico’s leading dairy producers, Data Protection and Compliance mandate: rather than treating data protection as a stand-alone documentation exercise, the engagement addressed how data-protection requirements should be incorporated during operational change. Grupo Excelsior instructed the firm to advise on the regularisation and restructuring of its operations, with a specific focus on data protection.

Arturo González de Araujo, Ernesto Zavala and Damián Ortega led the work. The team designed and implemented policies and training as practical components of the regularisation and restructuring process. Those measures were directed to embedding a data protection culture within Grupo Excelsior’s operations, while aligning the resulting arrangements with governance frameworks. The mandate therefore combined the operational task of restructuring and regularisation with concrete policy and training implementation, creating institutional measures intended to reduce regulatory exposure and place data protection within the client’s governance arrangements.
\item \textbf{Grupo Modelquipo Data Protection Strategy} (Client: Grupo Modelquipo | Value: N/A)\\ Embedding data protection within corporate governance, rather than isolating it as a discrete compliance exercise, is the central importance of the Grupo Modelquipo, Industrial equipment and machinery solutions provider, Data Protection Strategy. Grupo Modelquipo, improving compliance and integrating data protection into governance culture, engaged the firm in relation to a data protection framework developed as part of its corporate governance strategy. The mandate placed data protection directly within the client’s governance agenda, making the framework relevant not only to data-processing activity but also to the way governance considerations were addressed across the engagement.

The firm’s specific role was to advise on the data protection framework as part of Grupo Modelquipo’s corporate governance strategy. In performing that role, the team identified data processing risks and mitigated those risks through the framework and its integration into the corporate governance strategy. The work therefore connected the identification of data processing risks with their mitigation, while ensuring that data protection was addressed within the client’s governance culture rather than treated separately from it. The stated outcome was improved compliance and the integration of data protection into governance culture, reflecting the concrete governance-focused result recorded for Grupo Modelquipo.

Arturo González de Araujo, Ernesto Zavala and Damián Ortega led the matter. Their involvement encompassed the advice on the data protection framework, the identification and mitigation of data processing risks, and the integration of data protection into Grupo Modelquipo’s corporate governance strategy. The engagement demonstrates a governance-centred application of data protection advice: one directed toward improving compliance while embedding data protection into governance culture.
\item \textbf{Tiendas Chedraui Data Protection Assessment} (Client: Tiendas Chedraui | Value: N/A)\\ Tiendas Chedraui, Mexican supermarket and retail chain's expansion mandate demonstrates the practical importance of integrating data-protection assessment into the operational process for new store openings, where customer and operational data processing must be considered alongside the client’s growth plans. Tiendas Chedraui, supporting operational expansion while maintaining regulatory alignment, instructed the firm to advise on the legal assessment of data protection for new store openings and expansion.

The engagement focused on evaluating the risks connected with data processing in the context of planned new store openings and wider expansion. Rather than addressing data protection as an isolated compliance exercise, the firm’s role was to assess the relevant legal risks as Tiendas Chedraui pursued operational expansion, and to ensure data processing compliance in a manner aligned with those business activities. The assessment therefore connected the client’s proposed store-opening and expansion plans with the need for regulatory alignment in its data-processing arrangements.

Arturo González de Araujo and Damián Ortega led the matter. Their work involved the legal assessment of data protection, the evaluation of risks arising from the client’s data processing, and the compliance-focused analysis required for new store openings and expansion. The resulting advice supported Tiendas Chedraui’s operational expansion while maintaining regulatory alignment and ensuring data processing compliance.

\end{itemize}

\section*{IV. Strategic Evolution Path}
\begin{itemize}
\item Maintain current strategy.
\end{itemize}

\vfill
\centering
\small{This document is a structural assessment and does not guarantee ranking results.}

\end{document}